\chapter{结论}
\label{ch:conclusion}

\section{对四个研究问题的回答}
\label{sec:9-1}

\paragraph{RQ1：GPU 索引构建与容量。}
完整 JHQ 索引能够在单张 GPU 上构建和搜索，无须把原始数据库向量保存在设备上。训练和编码均在 GPU 执行，通过有界批处理与两遍构建控制临时存储。Stella 的 17.78 百万个向量若以 FP32 保存，需要 67.81\,GiB；JHQ-GPU 的常驻进程占用为 20.44\,GiB，包括上下文与工作区。索引冷态构建耗时 10.54--10.65\,s，在这张 32\,GB 设备上采样到的峰值为 22.58\,GiB（\cref{sec:7-7}）。

\paragraph{RQ2：面向表示结构的执行。}
JHQ 的编码结构支持小查找表、连续读取和并行候选选择。每张 256 项主码表拆成两张 16 项表后保持实数算术分数，并在 24 个配置中的 22 个成为最快的已测表形式。编码按字打包，使扫描相邻候选的线程读取连续内存；多个线程还可以协作精排一个候选。小 batch 下，将一个查询分到多个 block 能增加有效并行度：batch 为 8 时，调度相对采用协作精排的单 block 基线获得 $2.727\times$ 吞吐。分组选择在全部 55,296 个已查查询实例中，与 CPU 对 GPU 分数严格排序的结果一致；单 block 的差异仅出现在截断边界的并列候选上（\cref{sec:7-3,sec:7-6,sec:7-strict-selection}）。

\paragraph{RQ3：工作负载级精排校准。}
无标签查询能够选出可复用且满足已测质量目标的精排预算，但样本数量会影响结果。使用全部 500 条校准查询时，十二个配置的留出集召回率与已测预算中最高召回率的差距均在 $10^{-3}$ 以内，最大观测损失为 0.0006。小样本可能漏掉需要更多精排的查询：采用 128 条查询时，arXiv 在 128 probes 下有 64 次抽样中的 40 次超出目标，在 512 probes 下为 64 次中的 32 次。采用 128 查询校准时，OpenAI 直接复用测试相对固定预算对照在 4--9 个 batch 内收回开销。arXiv 选出的预算则需要花更多时间来达到质量目标（\cref{sec:7-5}）。

\paragraph{RQ4：跨系统搜索性能。}
主比较显示 JHQ-GPU 适合大 batch 吞吐场景。在召回率 0.97、batch 为 1,000 时，它在四个数据集上达到 CAGRA fp32 的 $1.32$--$1.75\times$ 吞吐；在已测网格中，扣除 CUDA 上下文后的净进程显存为 CAGRA 索引分配量的 0.29--0.39。它在 Vogue 和两个 OpenAI 数据集上领先已测方法，在 arXiv 上与 IVF-RaBitQ 接近持平；Vogue 的领先依赖分组扫描（\cref{sec:7-2}）。OpenAI 测量使用 TF32 乘法；复用相同训练模型的 FP32 检查仍领先 CAGRA fp32（\cref{sec:b-precision-anchors}）。这一优势受到 batch 大小限制：batch 为 8、召回率为 0.97 时，CAGRA fp32 在四个数据集上均更快（\cref{sec:7-smallbatch}）。

\section{贡献、局限与展望}
\label{sec:9-2}

当执行方式使编码访问成本降低、并行度匹配 batch、精排满足工作负载的质量需求时，压缩才能带来容量和速度收益。

证据覆盖一张 GPU、六个准备后的数据集、平方 L2 top-10 搜索和静态索引。历史 CPU 比较不能量化本文评测版本的加速比（\cref{sec:7-cpu}）。在校准池上与参考结果一致，也仍不能确定未来查询的质量。

由此可以继续研究两个方向。风险受控的校准应考虑工作负载漂移和预计复用次数，并与质量匹配的固定预算比较包含校准开销的总时间。从主机内存取回残差码有望增加容量；测量传输量、传输与计算的重叠以及尾延迟，可判断增加容量的代价是否值得（\cref{sec:8-4}）。
